\chapter{Imaginary-mode displacement (menu 30)}
\label{ch:menu30}
\menulabel{30}

\section{Purpose}

A converged optimisation that shows an imaginary frequency is not the
stationary point it claims to be.  The source's cure (the NIFREC workflow,
Tanaka \& Miyao; the software is MIT-licensed and its displacement stage was
read from the source) builds a vector from the imaginary modes --- the sum of
all of them by default, or the most negative mode alone --- normalizes it,
applies \code{coords + vec * disp}, and re-runs \code{opt freq} with
successively larger displacements (\code{base\_disp * (i + 1)}, default 0.1
Angstrom, up to five rounds) until every frequency is real.  The menu
generates the restart structures; the engine run stays with the user, and the
success criterion is the rerun's frequency check.

\section{What you need}

The frequency output that carries the imaginary mode (any \code{Opt}+\code{Freq}
output: the menu reads the last frequency block, the normal-mode vectors, and
the run's own final geometry and atomic masses) and the base input of that
run.  The base input's job keywords are kept verbatim --- a plain \code{Opt}
heals toward the minimum on each side of the mode while \code{OptTS} follows
the mode --- and only its coordinate block is replaced (with the \code{Freq}
token forced in, because the rerun's frequency check is the criterion).

\section{Walkthrough}

\lstinputlisting[style=fbkcode, firstline=35, lastline=51,
  title={The menu-30 example session (the F + H$_2$ fixture: the TS's
  $-90.48$~cm$^{-1}$ mode displaced $\pm 0.1$~\AA{} into the two restart
  inputs, each side carrying its largest single-atom move)}]
  {examples/expected/m30_imagdisp.txt}

\section{What you get}

\file{<input>.disp\_p.inp} and \file{<input>.disp\_m.inp} (the two signs of
the displacement --- an addition to the source, which grows one direction
0.1, 0.2, \ldots\ 0.5~\AA; both branches are then available in one round)
plus a report \file{<output>.imagdisp.fbk.md} recording the mode, the vector
choice and the amplitudes.  The ORCA-side conventions are measured: the
\code{NORMAL MODES} block prints mass-weighted vectors of unit Euclidean norm
(the Cartesian pattern is $v_i/\sqrt{m_i}$), and the masses come from the
run's own \code{(A.U.)} block, so the de-weighting uses the values ORCA's
Hessian was built with.

\section{Options worth knowing}

\begin{readbox}
\begin{itemize}
  \item \textbf{measured anchor} (F + H$_2$ fixture): the TS's
    $-90.48$~cm$^{-1}$ mode displaced $\pm 0.1$~\AA{} (largest single-atom
    move $0.055$~\AA) relaxes, both signs, to geometries with no imaginary
    frequencies at $-100.895757$~Eh (fixtures \file{fhh\_reopt.*}).  This
    mode is the near-linear bend of the shallow saddle, so the two signs
    land on the two mirror images of the same minimum --- the two-sided
    reaction separation needs an asymmetric-stretch mode, which this
    fixture does not carry (recorded rather than papered over);
  \item \textbf{if imaginary frequencies persist}: the source's schedule is
    the same vector with larger displacements (0.2, 0.3, \ldots\ up to
    0.5~\AA), or the other vector choice;
  \item \textbf{boundaries}: the displaced structure is a starting point,
    not a stationary point; the source's own limit applies --- the cure is
    probabilistic in the sense that a still-imaginary rerun means ``go
    larger or go elsewhere'', and the menu never claims the rerun will
    succeed.
\end{itemize}
\end{readbox}
